\documentclass[a4paper,12pt]{article}

\usepackage[utf8]{inputenc}
\usepackage[T1]{fontenc}
\usepackage[italian]{babel}
\usepackage{enumitem}
\usepackage{geometry}
\usepackage{hyperref}
\usepackage{booktabs}
\usepackage{graphicx}
\usepackage{longtable}

\geometry{margin=2.5cm}

\title{Documentazione del design  \\ \large Gestionale Locale per la Gestione del Magazzino}
\author{Gruppo 11}
\date{27 Maggio 2026}

\begin{document}

\maketitle

\section{Contesto e Obbiettivi}
Il progetto vuole creare uno strumento utilizzabile da un azienda per gestire dinamiche di magazzino , dinamiche di contabilita e di amministrazione.

\section{Non Obbiettivi}

Il sistema non intende affrontare la parte e-commerce cioè l'iserimento di ordini da parte del cliente finale.

\section{Proposta}

presenta la soluzione architetturale adottata, idealmente supportata da una singola figura riassuntiva.

\begin{itemize}
    \item \textbf{Front-end}: Go con Bubbletea per TUI
    \item \textbf{Back-end}: Abbiamo scelto GO per la velocita di calcolo e l'affidabilità nel campo della sicurezza
    \item \textbf{Database}: Utilizzaremo SQL lite per la semplicità di uso.
\end{itemize}

\section{Alternative}

elencano le opzioni valutate e i motivi per cui non sono state scelte.

\begin{itemize}
    \item \textbf{Poi} si farà
\end{itemize}

\section{Conseguenze}

analizzano vantaggi, svantaggi e rischi introdotti dalla soluzione proposta.

\begin{itemize}
    \item \textbf{Poi} si farà
\end{itemize}

\section{Piano di Adozione}

descrive come e quando la soluzione viene introdotta, se è prevista una migrazione o un’evoluzione graduale.

\begin{itemize}
    \item \textbf{Poi} si farà
\end{itemize}


\end{document}